\documentclass[aspectratio=169]{beamer}

% ============================================================
% Custom theme replicating the reference deck's look: sans-serif
% body, blue frame titles, hollow/filled triangle bullets, and a
% three-segment colored footline (institution | short title |
% date/page). Built from scratch (no \usetheme{...}) to match
% the reference slides exactly rather than approximate a stock
% Beamer theme.
% ============================================================

\usepackage[T1]{fontenc}
\usepackage{lmodern}
\usepackage[utf8]{inputenc}
\usepackage{amsmath, amssymb}
\usepackage{booktabs}
\usepackage{graphicx}
\usepackage[round]{natbib}

\setbeamertemplate{navigation symbols}{}

% ---- Colors -------------------------------------------------
\definecolor{slidetitleblue}{RGB}{60,90,190}
\definecolor{footdark}{RGB}{35,45,90}
\definecolor{footmid}{RGB}{70,95,175}

\setbeamercolor{structure}{fg=slidetitleblue}
\setbeamercolor{frametitle}{fg=slidetitleblue,bg=white}
\setbeamercolor{title}{fg=slidetitleblue}
\setbeamercolor{normal text}{fg=black,bg=white}
\setbeamercolor{footline author}{fg=white,bg=footdark}
\setbeamercolor{footline title}{fg=white,bg=footmid}
\setbeamercolor{footline date}{fg=white,bg=footdark}

\setbeamerfont{frametitle}{series=\mdseries,size=\Large}
\setbeamerfont{footline author}{size=\scriptsize}
\setbeamerfont{footline title}{size=\scriptsize}
\setbeamerfont{footline date}{size=\scriptsize}

% ---- Bullets: hollow triangle (level 1), filled triangle (level 2) --
\setbeamertemplate{itemize item}{\color{slidetitleblue}$\triangleright$}
\setbeamertemplate{itemize subitem}{\color{slidetitleblue}$\blacktriangleright$}
\setbeamertemplate{itemize subsubitem}{\color{slidetitleblue}$\bullet$}

\setbeamertemplate{frametitle}{%
  \vspace{1em}%
  \hspace{0.6em}\usebeamerfont{frametitle}\usebeamercolor[fg]{frametitle}\insertframetitle%
  \vspace{0pt}%
}
% Fixed, uniform gap between the title and the start of the body,
% applied after the frametitle regardless of what follows (itemize,
% enumerate, center, tabular, equation) so every slide starts its
% content at the same distance from the title.
\addtobeamertemplate{frametitle}{}{\vspace{0.35cm}}

% ---- Three-segment footline ---------------------------------
\newcommand{\shorttitletext}{Sovereign Spread Crises and the Sovereign-Bank Nexus}
\newcommand{\shortauthortext}{David Betsem}

\setbeamertemplate{footline}{%
  \leavevmode%
  \hbox{%
    \begin{beamercolorbox}[wd=.27\paperwidth,ht=2.4ex,dp=1.2ex,left,leftskip=0.6em]{footline author}%
      \usebeamerfont{footline author}\shortauthortext%
    \end{beamercolorbox}%
    \begin{beamercolorbox}[wd=.53\paperwidth,ht=2.4ex,dp=1.2ex,center]{footline title}%
      \usebeamerfont{footline title}\shorttitletext%
    \end{beamercolorbox}%
    \begin{beamercolorbox}[wd=.20\paperwidth,ht=2.4ex,dp=1.2ex,right,rightskip=0.6em]{footline date}%
      \usebeamerfont{footline date}\insertshortdate{} \hspace{0.4em} \insertframenumber{} / \inserttotalframenumber%
    \end{beamercolorbox}%
  }%
}

\setbeamertemplate{navigation symbols}{}
\setbeamertemplate{headline}{}

% ---- Title page (affiliations as superscripts, small disclaimer) ----
\title{\large Costs of Sovereign Debt Crises With and Without Default:\\
Effect of the Sovereign-Bank Nexus}
\author{\normalsize David Betsem A Djam}
\institute{%
  \footnotesize
  HEC Paris, Institut Polytechnique de Paris
}
\date{\normalsize 28\textsuperscript{th} September 2026}

\begin{document}

% ---------------- Title slide ----------------
{
\setbeamertemplate{footline}{}
\begin{frame}[plain]
  \vspace{0.8cm}
  \vfill
  \begin{center}
    \normalsize Master Thesis Defense\par
  \end{center}

  \begin{center}
    \usebeamercolor[fg]{title}{\Large Costs of Sovereign Debt Crises With and\\
    Without Default: Effect of the Sovereign-Bank Nexus\par}
  \end{center}
  \vspace{0.2cm}
  \begin{center}
    \large David Betsem\par
  \end{center}
  \vspace{0.6em}
  \begin{center}
    \small 28\textsuperscript{th} September 2026\par
  \end{center}
  \vfill
\end{frame}
}

% ---------------- Motivation ----------------
\begin{frame}{Motivation: sovereign debt crises do not always involve default}
  \begin{itemize}
    \item \textbf{Sovereign spread crises are more common than defaults}
    \begin{itemize}
      \item Episodes of severe bond-market distress (Mexico 1994/95, Thailand
      1997/98, Russia 2015) occurred without a missed payment or restructuring
      \item Elevated spreads reflect rising default risk even when no default
      is ultimately realized \citep{arellano2008default}
    \end{itemize}
    \item \textbf{The existing literature measures the cost of default}
    \begin{itemize}
      \item Output costs of default are estimated at 1--4 pp of GDP growth
      \item Rarely benchmarked against equally severe episodes that did not
      end in default
    \end{itemize}
    \item \textbf{This paper compares both types of crisis directly}
    \begin{itemize}
      \item Do spread crises without default carry a materially different
      cost than default-linked crises?
      \item Does the pre-crisis sovereign-bank nexus shape that difference?
    \end{itemize}
  \end{itemize}
\end{frame}

% ---------------- This paper ----------------
\begin{frame}{This paper}
  \begin{itemize}
    \item \textbf{Research question:} How do GDP, bank credit, and investment
    respond to sovereign spread crises with and without default, and does the
    pre-crisis sovereign-bank nexus amplify that response?
    \item \textbf{Data:} JP Morgan EMBIG spread database, 1993--2023
    \begin{itemize}
      \item 61 spread crisis episodes, 21 default-linked, 40 non-default
    \end{itemize}
    \item \textbf{Method:} Local projection with Augmented Inverse Probability
    Weighting (AIPW), following \citet{jorda2016time}
    \item \textbf{Identification:} Comparison across crisis type and
    pre-crisis sovereign-bank exposure (median split on bank claims on
    government / assets)
  \end{itemize}
\end{frame}

% ---------------- Related literature ----------------
\begin{frame}{Related literature}
  \small
  \begin{itemize}
    \item \textbf{Debt crises without default}: \citet{pescatori2007adequately}, \citet{broner2013why}, \citet{aguiar2016quantitative}, \citet{krishnamurthy2017credit}, \citet{mitchener2021sovereign}
    \item \textbf{Fluctuations in country spreads}: \citet{uribeyue2006country}, \citet{arellano2008default}, \citet{mendozayue2012general}, \citet{bocola2016passthrough}, \citet{bahaj2020sovereign}
    \item \textbf{Output costs of default}: \citet{borensztein2009costs}, \citet{furceri2012costly}, \citet{trebesch2017output}, \citet{kuvshinov2019sovereigns}, \citet{levyyeyati2011elusive}, \citet{asonuma2024costs}
    \item \textbf{Sovereign-bank nexus}: \citet{gennaioli2014sovereign}, \citet{gennaioli2018banks}, \citet{acharya2015greatest}, \citet{dellariccia2018managing}, \citet{broner2014sovereign}, \citet{angelini2014negative}
  \end{itemize}
  \vspace{0.6em}
  \noindent This paper contributes by jointly comparing default-linked
  and non-default crises and their pre-crisis sovereign-bank exposure.
\end{frame}

% ---------------- Preview of results ----------------
\begin{frame}{Preview of results}
  \small
  \noindent All figures: AIPW local projection estimates, own pre-crisis
  path = 0.\par
  \vspace{0.4em}
  \begin{itemize}
    \item \textbf{Crisis type is predictable} (at onset): larger spread
    shocks, more trade openness, sharper depreciation $\to$ default;
    deeper bank credit, higher inflation $\to$ non-default
    \item \textbf{GDP} ($h=1$--5): $-$1--2 pp (non-default) vs.\
    $-$4--9 pp (default)
    \item \textbf{Bank credit} ($h=1$--5): +7 pp by $h=5$
    (non-default) vs.\ $-$40+ pp (default)
    \item \textbf{Investment, FDI, lending rates} ($h=1$--5): no
    significant difference by crisis type
    \item \textbf{Claims on government} ($h=1$): +7 pp vs.\ $<$1 pp --
    only outcome robust under both tests
    \item \textbf{Nexus} ($h=1$): matters for non-default (0.75 vs.\
    $-$2.87 pp, sig.), not for default (18 vs.\ 8 pp, not sig.)
  \end{itemize}
\end{frame}

% ---------------- Roadmap ----------------
\begin{frame}{Roadmap}
  \begin{enumerate}
    \item \textbf{Data and crisis classification}
    \item Estimation strategy and main results
    \item Role of the sovereign-bank nexus
    \item Conclusion
  \end{enumerate}
\end{frame}

% ---------------- Section divider example ----------------
\begin{frame}{Data sources}
  \begin{itemize}
    \item \textbf{JP Morgan EMBIG Global}
    \begin{itemize}
      \item Sovereign bond spread series used to identify crisis onsets
    \end{itemize}
    \item \textbf{IMF WEO, IMF FSI, World Bank IDS}
    \begin{itemize}
      \item GDP, investment, bank credit, bank claims on government, public
      debt, trade openness
    \end{itemize}
    \item \textbf{\citet{asonuma2016sovereign}}
    \begin{itemize}
      \item Sovereign restructuring database, used to classify episodes as
      default-linked or non-default
    \end{itemize}
    \item \textbf{\citet{laeven2013systemic,laeven2026systemic}}
    \begin{itemize}
      \item Systemic banking crisis dummy
    \end{itemize}
  \end{itemize}
\end{frame}

% ---------------- Why a first-stage probit? ----------------
\begin{frame}{Why a first-stage probit?}
  \small
  \begin{itemize}
    \item Crises are not randomly assigned: countries entering a
    default-linked or non-default spread crisis are systematically
    different on observables, and those same pre-crisis conditions can
    also shape the post-crisis path
    \item A simple comparison of means, or plain OLS, would conflate
    this selection with the effect of the crisis itself
    \item Solution: estimate each country-year's ex-ante probability of
    entering each crisis type, and use it to reweight observations
    (AIPW) so the estimated effect is drawn from a distribution closer
    to what would prevail absent this predictable selection
    \citep{altonji2005selection}
  \end{itemize}
\end{frame}

% ---------------- Main Specification: Probit model ----------------
\begin{frame}{Main Specification: Probit model}
  \footnotesize

  We estimate each country-year's ex-ante probability of entering a
  non-default or default-linked spread crisis, and derive the
  predicted probability $\widehat{p}_{i,t+1}^S$
  \[
  \Pr\!\left(D_{i,t+1}^S = 1 \mid X_{i,t}, Z_{i,t}\right) = \Phi\!\left(X_{i,t}\gamma^{S,X}, Z_{i,t}\gamma^{S,Z}\right)
  \]
  \vspace{-0.2em}
  \begin{itemize}
    \item $D_{i,t+1}^S \in \{0,1\}$: dummy equal to 1 if country $i$
    enters a crisis of type $S$ in year $t+1$
    \item $S \in \{$non-default, default-linked$\}$: two probits
    estimated separately, one per crisis type
    \item $\Phi(\cdot)$: standard normal CDF (probit link function)
    \item $X_{i,t}$: 8 baseline controls (GDP growth, public
    debt-to-GDP, banking crisis dummy, government expenditure-to-GDP,
    trade openness, bank credit-to-GDP, inflation, nominal exchange
    rate)
    \item $Z_{i,t}$: 3 predictors -- US federal funds rate, a
    contagion measure, years since the country's last default-linked
    onset
    \item $\gamma^{S,X}$, $\gamma^{S,Z}$: estimated coefficients on
    $X_{i,t}$ and $Z_{i,t}$, respectively, for crisis type $S$
  \end{itemize}
\end{frame}

% ---------------- Probit: results ----------------
\begin{frame}{First stage: what predicts crisis onset and default?}
  \small
  \begin{itemize}
    \item \textbf{Onset vs.\ tranquil years} (Table~1): US federal
    funds rate and years since last default-linked onset predict both
    crisis types; GDP growth, bank credit-to-GDP, inflation and the
    exchange rate predict default-linked onset specifically, while no
    domestic control significantly predicts non-default onset
    \item AUROC: 0.77 (non-default) and 0.89 (default-linked) with
    predictors, vs.\ 0.65 and 0.84 with controls only
    \item \textbf{Robustness -- default vs.\ non-default among onsets
    only:} peak EMBIG spread at onset significantly predicts a
    default-linked resolution, even conditional on the same controls
    \begin{itemize}
      \item Trade openness and currency depreciation also raise the
      likelihood of default-linked resolution; deeper bank credit and
      higher inflation lower it
    \end{itemize}
  \end{itemize}
\end{frame}

% ---------------- Main Specification: outcome model ----------------
\begin{frame}{Main Specification: outcome model}
  \footnotesize

  For each crisis type $S$ and horizon $h$, we estimate a single
  outcome regression, weighted by the Stage 1 propensity score
  (IPWRA: inverse-probability-weighted regression adjustment)
  \[
  y_{i,t+h} - y_{i,t} = \alpha_i^{S,h} + \beta^{S,h} D_{i,t+1}^{S} + X_{i,t}\beta^{S,h} + \mu_{i,t+h}^{S}
  \]
  \vspace{-0.2em}
  \begin{itemize}
    \item $y_{i,t+h} - y_{i,t}$: cumulative change in the outcome $y$
    from crisis onset ($t$) to horizon $h$
    \item $D_{i,t+1}^{S}$: same onset dummy as in the probit, for
    crisis type $S$ -- entered as a regressor, not split into two
    samples
    \item $\alpha_i^{S,h}$: country fixed effects, specific to crisis
    type $S$ and horizon $h$
    \item $X_{i,t}$: same 8 baseline controls as the probit stage
    \item $\beta^{S,h}$: coefficient on $D$ and on $X_{i,t}$
    \item $\mu_{i,t+h}^{S}$: error term
    \item One regression, weighted by
    $w_{i,t} = D_{i,t+1}^S/\widehat{p}_{i,t+1}^S + (1-D_{i,t+1}^S)/(1-\widehat{p}_{i,t+1}^S)$
    (from Stage 1)
    \item $\to$ predict twice from the same fitted regression, once
    setting $D=1$ and once $D=0$, to get $\widehat{m}_1^{S,h}$ and
    $\widehat{m}_0^{S,h}$ -- inputs to the AIPW estimator, not the
    final effect
  \end{itemize}
\end{frame}

% ---------------- Main Specification: combining into AIPW ----------------
\begin{frame}{Main Specification: the AIPW estimator}
  \footnotesize

  We combine the outcome model ($\widehat{m}_1^{S,h}$,
  $\widehat{m}_0^{S,h}$) and the probit ($\widehat{p}_{i,t+1}^S$) into
  the AIPW estimator, following \citet{jorda2016time}:
  \[
  \Lambda_h^S = \frac{1}{N}\sum_{i}\sum_{t}\Biggl\{
  \underbrace{\left[\frac{D_{i,t+1}^S(y_{i,t+h}-y_{i,t})}{\widehat{p}_{i,t+1}^S}
    - \frac{(1-D_{i,t+1}^S)(y_{i,t+h}-y_{i,t})}{1-\widehat{p}_{i,t+1}^S}\right]}_{\text{IPW term}}
  - \underbrace{\frac{D_{i,t+1}^S - \widehat{p}_{i,t+1}^S}{\widehat{p}_{i,t+1}^S(1-\widehat{p}_{i,t+1}^S)}
    \left[(1-\widehat{p}_{i,t+1}^S)\widehat{m}_1^{S,h} + \widehat{p}_{i,t+1}^S\widehat{m}_0^{S,h}\right]}_{\text{augmentation term}}
  \Biggr\}
  \]
  \vspace{-0.2em}
  \begin{itemize}
    \item $\Lambda_h^S$: the AIPW estimate -- the effect of crisis type
    $S$ on the outcome at horizon $h$, our object of interest
    \item \textbf{Not} simply reweighting observations and running the
    local projection: the augmentation term is required
    \item Doubly robust: $\Lambda_h^S$ stays consistent if either the
    probit \emph{or} the outcome model is correctly specified, not both
    \item Differences between non-default and default-linked estimates
    are assessed with a paired bootstrap test (governing) and the
    \citet{clogg1995statistical} test (companion, assumes independence)
  \end{itemize}
  \vspace{0.3em}
  \begin{center}
    \small
    Outcome model ($\widehat{m}_1,\widehat{m}_0$) $+$ Probit
    ($\widehat{p}$) $\;\to\;$ AIPW estimator $\;\to\;$ effect by
    horizon $h$
  \end{center}
\end{frame}

% ---------------- AIPW LP: results, GDP ----------------
\begin{frame}{Result: GDP}
  \begin{itemize}
    \item Non-default crises: GDP declines modestly, by around
    $-$1--2 pp over $h=1$--2, before returning towards its pre-crisis
    path
    \item Default-linked crises: GDP falls by $-$4--9 pp, with the
    decline remaining significant through $h=5$
    \item The gap between the two crisis types is significant under
    the bootstrap test
  \end{itemize}
\end{frame}

% ---------------- AIPW LP: results, bank credit ----------------
\begin{frame}{Result: bank credit to the private sector}
  \begin{itemize}
    \item Non-default crises: bank credit turns significantly positive
    by $h=5$, reaching nearly +7 pp above its pre-crisis path
    \item Default-linked crises: bank credit falls to more than
    $-$40 pp below its pre-crisis path, a contraction that persists
    through the end of the horizon
    \item A similar contrast to GDP, but more delayed in onset
  \end{itemize}
\end{frame}

% ---------------- AIPW LP: results, investment/FDI/lending ----------------
\begin{frame}{Result: investment, FDI, and lending rates}
  \begin{itemize}
    \item All three move significantly following at least one crisis
    type
    \item But the difference between non-default and default-linked
    crises is never significant under the bootstrap test, at any
    horizon
    \item A genuine null result -- not merely a smaller-magnitude
    difference
  \end{itemize}
\end{frame}

% ---------------- AIPW LP: results, claims on government ----------------
\begin{frame}{Result: bank claims on government}
  \begin{itemize}
    \item Default-linked crises: claims on government rise
    significantly, by approximately +7 pp at $h=1$
    \item Non-default crises: claims on government decline by less
    than 1 pp at $h=1$
    \item This difference is significant under the bootstrap test --
    the only outcome robust under both the bootstrap and the Clogg test
  \end{itemize}
\end{frame}

% ---------------- Nexus specification ----------------
\begin{frame}{Sovereign-bank nexus: specification}
  \begin{itemize}
    \item \textbf{Measure:} banks' claims on government as a share of
    total bank assets, lagged one year before crisis onset; median
    split (14.96\%) into high- and low-nexus groups, 25 onsets each
    \item \textbf{Specification:} interact the crisis-type dummy with
    the high-nexus dummy inside the same local-projection framework
  \end{itemize}
  \vspace{0.3em}
  \begin{center}
    \footnotesize
    \[
    y_{i,t+h} - y_{i,t} = \delta_i^{S,h}
      + \Gamma^{HN,h} D_{i,t}^{HN} D_{i,t+1}^{S}
      + \Gamma^{LN,h} \left(1 - D_{i,t}^{HN}\right) D_{i,t+1}^{S}
      + X_{i,t}\Gamma^{S,h} + e_{i,t+h}^{S}
    \]
  \end{center}
  \vspace{0.3em}
  \begin{itemize}
    \item $\Gamma^{HN,h}$, $\Gamma^{LN,h}$: effect of a crisis of type
    $S$ when pre-crisis sovereign exposure is high or low,
    respectively; differences assessed via paired bootstrap and the
    \citet{clogg1995statistical} test
  \end{itemize}
\end{frame}

% ---------------- Nexus results ----------------
\begin{frame}{Sovereign-bank nexus: result}
  \begin{itemize}
    \item \textbf{Non-default crises} ($h=1$): GDP changes by
    approximately $+$0.75 pp in the high-nexus group, vs.\ a decline of
    around $-$2.87 pp in the low-nexus group -- significant under the
    bootstrap test
    \item \textbf{Default-linked crises} ($h=1$): GDP declines by
    approximately $-$18 pp in the high-nexus group, vs.\ around
    $-$8 pp in the low-nexus group -- not significant under the
    bootstrap test, at any horizon
    \item No systematic evidence that a stronger pre-crisis nexus
    amplifies the output costs of default-linked crises; it is
    associated with a smaller GDP decline on impact for non-default
    crises
  \end{itemize}
\end{frame}

% ---------------- Conclusion ----------------
\begin{frame}{Conclusion}
  \begin{itemize}
    \item Sovereign spread crises without default carry real but modest
    output costs, far smaller than default-linked crises
    \item The sovereign-bank nexus amplifies the recovery in non-default
    crises, consistent with a bank-lending channel that is only active when
    the sovereign avoids default
    \item No significant nexus effect within default-linked crises, though
    the point estimates are consistently more negative for high-nexus
    countries -- a null result given the limited sample of default-linked
    episodes, not evidence of no effect
  \end{itemize}
\end{frame}

\begin{frame}[allowframebreaks]{References}
  \footnotesize
  \bibliographystyle{apalike}
  \bibliography{references}
\end{frame}

\end{document}
